%%
%% Trace Data Quality for LLM-Agent Observability:
%% A Controlled Study of Schema Loss and Diagnostic Degradation in
%% Telemetry Backends
%%
%% Reconstructed LaTeX source of the submitted IEEE BigData 2026 paper
%% (submitted/BigD260_4648.pdf). The original source was lost; this file
%% reproduces the submitted text, tables, and references verbatim so that
%% later revisions can be diffed against it. Content is intentionally NOT
%% edited, corrected, or anonymized.
%%
%% Build: tectonic main.tex
%% The original was produced by pdfTeX (TeX Live 2025), where OT1 is the
%% default font encoding. The explicit [OT1]{fontenc} below keeps the same
%% Times/Courier OT1 fonts when the file is compiled with XeTeX-based
%% tectonic; it is a no-op under pdfLaTeX.
%%

\documentclass[conference]{IEEEtran}
\IEEEoverridecommandlockouts

\usepackage[OT1]{fontenc}
\usepackage{cite}
\usepackage{booktabs}
\usepackage{url}
\usepackage{enumitem}
\usepackage{amsmath}
\usepackage{graphicx}
\usepackage{microtype}
\usepackage{balance}

\begin{document}

\title{Trace Data Quality for LLM-Agent Observability:
A Controlled Study of Schema Loss and Diagnostic
Degradation in Telemetry Backends}

\author{\IEEEauthorblockN{Krishna Chaitanya Balusu}
\IEEEauthorblockA{Independent Researcher\\
San Francisco, CA, USA\\
krishnabkc15@gmail.com}}

\maketitle

\begin{abstract}
LLM-agent systems generate execution traces that look superficially
similar to ordinary distributed traces but carry different diagnostic
semantics: planning steps, reasoning loops, guardrail decisions,
delegation, memory accesses, retrieval, tool calls, and LLM invocations
need to remain distinguishable for the structural diagnostics studied
here after collection, export, and downstream analysis. Existing
observability schemas preserve some LLM and tool events, but often
collapse or remove agent-orchestration events. This paper treats
LLM-agent observability as a big-data quality problem. We define a
trace-data quality model covering schema coverage, semantic
preservation, joinability, payload-redacted structural sufficiency, and
diagnostic-proxy sufficiency. Using existing open AgentTelemetry
artifacts as measurement data and independent $\tau$-bench
trajectories, we analyze 2,148 traces containing 46,267 events exported
as OpenTelemetry spans: 112 SWE-bench traces, a separate 56-trace
end-to-end/topology suite, and 1,980 $\tau$-bench tool-agent
trajectories. We project those events into seven schemas, export each
projection through the OpenTelemetry SDK into a SQLite trace backend,
evaluate eight structural visibility queries, and run field ablations
over the reference schema. We also validate an OTel Collector
receiver/exporter round trip for the custom-attribute projection.
LLM-focused projections preserve tool, token, and retrieval queries, but
retain only 33.3--55.6\% of semantic event kinds and answer
37.5--50.0\% of positive structural visibility query families overall.
They answer 0.0\% of planning, reasoning-loop, memory, and delegation
queries. Removing only parent identifiers drops query recall to 87.5\%
by breaking tool provenance; removing event-kind fields drops query
recall to 0.0\%. A payload-redacted structural projection retains
100.0\% recall for measured structural visibility queries while removing
prompt/completion/tool payload content. An OTel+custom-attributes
projection shows the same structural recall is possible in generic span
storage when non-content diagnostic fields are retained. The result is a
reproducible controlled field-retention study of where LLM-agent trace
pipelines lose structural answerability and which schema properties the
measured queries require.
\end{abstract}

\begin{IEEEkeywords}
big data, observability, LLM agents, data quality, tracing,
OpenTelemetry, diagnostic analytics
\end{IEEEkeywords}

%% =====================================================================
\section{Introduction}
\label{sec:introduction}

Distributed tracing was designed to make large software systems
observable by converting executions into structured event streams. That
assumption becomes fragile for LLM-agent systems. An agent run is not
just an RPC path or a sequence of service calls; it is a data-rich
process involving planning, reasoning, retrieval, tool selection,
guardrail checks, delegation, memory reads and writes, and repeated LLM
invocations. If a telemetry pipeline preserves only generic spans,
tokens, and tool calls, downstream analytics cannot distinguish a
healthy multi-step plan from a reasoning loop, a blocked guardrail from
an ordinary application branch, or a stale memory read from a normal
lookup.

This is a data-quality problem. Agent traces are collected at high
volume, transformed across incompatible schemas, exported into tracing
backends, and queried by incident-response and analytics workflows. The
core question is not only whether traces are present, but whether they
retain the semantic fields required for later joins, aggregation,
privacy filtering, and diagnosis. In big-data terms, LLM-agent
observability raises a veracity and interoperability problem: the
pipeline may contain many spans while losing the information that makes
those spans analytically meaningful.

OpenTelemetry (OTel)~\cite{otel_spec} and its GenAI semantic conventions
provide a foundation for LLM invocation and tool execution
telemetry~\cite{otel_genai}. LLM-observability systems such as
OpenLLMetry~\cite{openllmetry}, LangSmith~\cite{langsmith},
Langfuse~\cite{langfuse}, and OpenInference~\cite{openinference} extend
that ecosystem. However, the agent layer adds orchestration semantics
that are not consistently preserved across schemas: planning, reasoning,
guardrails, delegation, and memory. The practical effect is schema loss:
trace rows continue to exist, but columns and event categories needed
for diagnosis are missing, collapsed, or no longer joinable.

This paper asks how to measure that loss. We use the public
\textsc{AgentTelemetry} artifacts from AIware
2026~\cite{agenttelemetry_aiware,agenttelemetry_artifact} as data
sources, not as the novelty of this paper. The prior work introduced a
benchmark, span taxonomy, fault classes, and toolkit. Here, the
contribution is different: a data-engineering analysis of trace-data
quality and diagnostic degradation across telemetry schemas. In
particular, we ask:

\begin{itemize}[nosep,leftmargin=*]
\item \textbf{RQ1}: What trace data are lost when LLM-agent executions
  are mapped into generic or LLM-only telemetry schemas?
\item \textbf{RQ2}: How does schema loss affect downstream diagnostic
  analytics such as loop detection, root-cause localization, and spans
  to diagnosis?
\item \textbf{RQ3}: Which schema properties are sufficient for
  payload-redacted structural queryability, and which calibration risks
  remain even when agent semantics are preserved?
\end{itemize}

The analysis combines three pieces of evidence. First, we perform a
schema-interoperability study over 2,148 traces containing 46,267 events
exported as OpenTelemetry spans: 112 SWE-bench~\cite{swebench} traces, a
separate 56-trace end-to-end/topology suite, and 1,980
$\tau$-bench~\cite{taubench} tool-agent trajectories. The $\tau$-bench
workload provides independent task outcomes and expected tool-action
labels, so two query families are evaluated against external task
labels rather than agent event-kind labels. Second, we run eight
downstream structural visibility queries against seven schema
projections, measuring which answers survive schema transformation after
export through the OpenTelemetry SDK into a SQLite trace backend. Third,
we run field ablations that remove one class of fields at a
time---event kind, parent identifiers, timestamps, memory fields,
guardrail fields, token/cost fields, and content payloads. A 3,780-row
controlled fault matrix is used only as validation that the schema
distinction aligns with fault-family detection.

The results are concrete but deliberately scoped. LLM-focused
projections preserve tool, token, and retrieval queries, but retain only
33.3--55.6\% of semantic event kinds and answer only 37.5--50.0\% of
positive structural visibility query families overall. They answer
0.0\% of planning, reasoning-loop, memory, and delegation queries
because the required orchestration fields are not present. Field
ablations show the failure modes are not all-or-nothing: dropping only
parent identifiers breaks tool provenance, dropping only timestamps
breaks reasoning-loop detection, and dropping only memory or guardrail
attributes breaks the corresponding audits. Conversely, a
payload-redacted structural projection preserves all measured
structural visibility answers without retaining prompt, completion,
reasoning-chain text, tool payload, or retrieved query content.

\textbf{Contributions.}
\begin{itemize}[nosep,leftmargin=*]
\item A trace-data quality model for LLM-agent observability covering
  schema coverage, semantic preservation, joinability, payload-redacted
  structural sufficiency, and diagnostic sufficiency.
\item A reproducible schema-interoperability study over 2,148 traces and
  46,267 source events materialized as 740,272 backend span rows across
  projections, separating row retention, semantic-kind coverage,
  required-attribute preservation, join-key preservation, and query
  answerability.
\item Empirical evidence that schema loss degrades downstream structural
  diagnostic answerability: LLM-focused projections answer
  37.5--50.0\% of positive structural visibility query families overall
  and 0.0\% of the measured planning, reasoning-loop, memory, and
  delegation families.
\item A separate $\tau$-bench action-visibility control that measures
  expected-action and argument-shape visibility, then shows that a
  missing-action proxy is weak against task pass/fail baselines.
\item A custom-attribute projection showing that generic OTel-style
  storage can preserve structural queryability when non-content
  diagnostic fields are retained.
\item A field-ablation study showing which individual field classes are
  load-bearing for reasoning-loop, guardrail, memory, delegation,
  tool-provenance, token-accounting, and retrieval-grounding queries.
\item Generated tables and scripts that rederive the paper's metrics
  from sanitized derived artifacts without running new LLM experiments.
\end{itemize}

%% =====================================================================
\section{Background and Boundary}
\label{sec:background}

\subsection{Trace data in LLM-agent systems}

Classical distributed tracing systems such as Dapper, X-Trace, and
Magpie model request execution as structured event streams across
services~\cite{dapper,xtrace,magpie}. A span typically has a name,
timestamp, duration, parent, trace identifier, service attributes, and
semantic-convention fields. That model is still useful for LLM-agent
systems, but it is not enough. A single agent trace may contain an LLM
call whose output is a plan, a retrieval call that grounds the plan, a
tool call selected by a reasoning step, a guardrail check, and memory
updates used by later turns. These are not interchangeable internal
spans. They define the analytical schema of the trace.

OTel GenAI conventions standardize parts of this schema, especially LLM
invocation and tool execution~\cite{otel_genai}. OpenInference and
OpenLLMetry-style systems extend the space with LLM and retrieval
observability~\cite{openllmetry,openinference}. Yet agent failures
reported in recent work include reasoning loops, delegation cycles,
planning failures, memory corruption, and guardrail bypass
events~\cite{mast,agentops,agdebugger}. These failures require fields
that are not always retained by LLM-only tracing.

\subsection{Boundary with prior AgentTelemetry work}

This paper deliberately does not claim the novelty of the
\textsc{AgentTelemetry} span taxonomy, fault classes, or SDK. Those
belong to the AIware 2026 paper~\cite{agenttelemetry_aiware}. We cite
that work as the source of the measurement artifacts and use its result
files as input data. We also do not claim a new adapter architecture or
runtime middleware layer. The contribution here is the data-quality
analysis: what happens to structural diagnostic utility when trace
schemas preserve or drop particular semantic fields.

This boundary matters for both scientific clarity and publication
ethics. A benchmark/toolkit paper asks whether a set of faults and spans
can be measured. A middleware paper asks how instrumentation can be
attached to heterogeneous frameworks. This BigData paper asks a
different question: once traces are produced, which schema properties
are retained by the big-data pipeline for downstream analytics to remain
valid?

%% =====================================================================
\section{Trace-Data Quality Model}
\label{sec:model}

We model an agent trace as a set of spans $T=\{s_1,\ldots,s_n\}$ with
parent-child edges, timestamps, span kinds, and key-value attributes. A
telemetry pipeline applies a schema mapping $m$ that may retain spans,
remove spans, collapse kinds, remove attributes, or preserve only
structural fields. A data-quality evaluation must therefore distinguish
\emph{row retention} from \emph{semantic preservation}: a trace may
retain many rows while losing the columns needed for diagnosis.

\subsection{Metrics}

\textbf{Schema coverage} measures the fraction of observed semantic
event kinds that remain visible after mapping:
\[
C(m,T) = \frac{|K(m(T)) \cap K(T)|}{|K(T)|},
\]
where $K(T)$ is the set of event kinds present in the original trace.

\textbf{Span retention} measures the fraction of spans that remain
analytically typed after mapping:
\[
R(m,T) = \frac{|\{s \in T : \mathrm{kind}(m(s)) \neq \bot\}|}{|T|}.
\]
This differs from raw export count: a backend may store a span but strip
the attribute that makes it usable for agent analytics.

\textbf{Semantic preservation} is a task-specific metric. For a
diagnostic query $q$ (for example, ``find reasoning loops''), it asks
whether the mapped trace contains the fields required to answer $q$. We
use binary preservation for queries whose required fields are known:
planning visibility, reasoning-loop visibility, guardrail event
visibility, memory-operation visibility, and tool-error visibility.

\textbf{Joinability} measures whether events can be linked by trace
identifier, parent-child relation, agent identifier, or task/session
identifier. In this study we measure joinability qualitatively through
diagnostic outcomes rather than a separate numeric score: if schema loss
forces a diagnostic algorithm to inspect more candidate spans,
joinability has degraded for that task.

\textbf{Payload-redacted structural sufficiency} asks whether structural
fields are enough for a diagnostic task without prompt, completion, tool
payload, reasoning-chain, or retrieved-query content. This metric does
not claim that structural fields are risk-free; it asks whether a
pipeline can remove high-sensitivity payload content while preserving
planning, reasoning, memory, guardrail, and provenance structure.

\textbf{Diagnostic sufficiency} measures downstream utility: root-cause
candidate set size, spans to diagnosis, signal-to-noise ratio, and
fault-detection rate for fault families. These are not schema metrics by
themselves; they are the effects that make schema quality operationally
meaningful.

Table~\ref{tab:queries} makes the model operational. Each row is a
downstream analytic query that a data platform might run over agent
traces. The required fields are not exotic: event kind, parent
relations, agent/task identifiers, and a few diagnostic attributes. The
problem is that these fields cross schema layers. Some are produced by
framework adapters, some are standardized by semantic conventions, some
are preserved or dropped by collector processors, and some are redacted
by privacy filters. A data-quality model must therefore follow the trace
across the pipeline rather than treating export success as the end of
the problem.

\begin{table*}[t]
\centering
\scriptsize
\caption{Diagnostic queries and the trace-data fields they require. The
table illustrates why row retention alone is not enough: a pipeline can
export many spans while deleting the fields required by the query.}
\label{tab:queries}
\begin{tabular}{@{}p{0.19\textwidth}p{0.34\textwidth}p{0.39\textwidth}@{}}
\toprule
\textbf{Diagnostic query} & \textbf{Minimum trace-data requirements} & \textbf{Failure mode under schema loss} \\
\midrule
Reasoning-loop detection & Reasoning event kind, agent/task identifier, iteration or parent-child order, token/cost fields when available & Repeated reasoning becomes a sequence of generic internal spans or LLM calls; loop detection falls back to brittle name matching or cannot run. \\
\addlinespace
Guardrail audit & Guardrail event kind, policy/check identifier, decision outcome, relation to the triggering LLM/tool span & Blocked or bypassed safety events disappear from trace analytics; audit jobs cannot distinguish absent checks from passing checks. \\
\addlinespace
Planning failure analysis & Planning event kind, plan representation or step count, parent agent, downstream tool/retrieval children & The analyst cannot connect an invalid action sequence to the plan that generated it; root-cause localization expands to many candidate spans. \\
\addlinespace
Memory-corruption analysis & Memory event kind, operation, key or namespace, agent/session identity, relation to later reasoning/LLM spans & Stale or corrupted memory reads are indistinguishable from ordinary application lookups or are omitted entirely. \\
\addlinespace
Tool-decision provenance & Tool-call event, parent LLM/reasoning relation, selected tool name, status/error, nearby retrieved context & Wrong-tool and tool-error analysis can see the tool call but not the reasoning path that selected it. \\
\bottomrule
\end{tabular}
\end{table*}

\begin{table*}[t]
\centering
\scriptsize
\caption{Sanitized trace corpora used by the backend schema-quality
study. The $\tau$-bench workload is independent of the AgentTelemetry
artifacts and contributes task outcomes plus expected tool-action
labels.}
\label{tab:corpora}
\begin{tabular}{@{}lrrp{0.62\textwidth}@{}}
\toprule
Corpus & Traces & Exported events & Event-kind distribution \\
\midrule
SWE-bench traces & 112 & 3060 & AGENT:112, GUARD\_RAIL:34, LLM\_CALL:861, MEMORY:224, PLANNING:112, REASONING:914, RETRIEVAL:669, TOOL\_CALL:134 \\
E2E/topology traces & 56 & 626 & AGENT:176, DELEGATION:109, GUARD\_RAIL:46, LLM\_CALL:150, MEMORY:33, PLANNING:50, REASONING:38, RETRIEVAL:16, TOOL\_CALL:8 \\
$\tau$-bench trajectories & 1980 & 42581 & AGENT:1980, LLM\_CALL:26316, TOOL\_CALL:14285 \\
\bottomrule
\end{tabular}
\end{table*}

%% =====================================================================
\section{Methodology}
\label{sec:method}

\subsection{Data sources}

We use existing open result artifacts from the \textsc{AgentTelemetry}
artifact package~\cite{agenttelemetry_artifact} and independent
historical $\tau$-bench trajectories~\cite{taubench}. No new LLM calls
are required for the measurements in this paper; the new work is the
schema-projection, backend-export, query-answerability, and
field-ablation analysis over sanitized derived events.

\textbf{SWE-bench trace corpus.} The first corpus contains 112 SWE-bench
agent traces with 3,060 spans. These traces include planning, memory,
reasoning, LLM, retrieval, agent, tool, and guardrail events. We use
this corpus to measure schema loss and diagnostic-proxy changes under
degraded schema views.

\textbf{E2E/topology trace suite.} The second corpus contains 56 traces
with 626 spans from end-to-end and multi-agent topology runs. Unlike the
SWE-bench corpus, this suite includes explicit delegation events and
broader topology structure. It gives the schema study an independent
workload slice within the artifact package rather than only the
controlled fault matrix.

\textbf{$\tau$-bench trajectory corpus.} The third corpus contains 1,980
historical $\tau$-bench trajectories with 42,581 derived events. Each
trajectory has a task outcome and expected tool-action labels supplied
by the benchmark. We convert each trajectory into a sanitized event
trace containing an agent task event, assistant-call events, and
tool-call events. The conversion records tool names, argument-shape
counts, sequence order, and task outcome, but it does not write user
messages, system prompts, tool outputs, or raw argument values into the
paper artifacts. This workload is independent of the
\textsc{AgentTelemetry} artifact and supports an external tool-sequence
visibility control.

\textbf{Controlled fault matrix.} The fourth result set is a 3,780-row
fault-injection result matrix containing fourteen non-empty fault
classes plus no-fault controls, six telemetry conditions, seven
frameworks, and six LLMs. We use this matrix as validation for the
fault-family interpretation, not as the center of the new contribution.

\subsection{Schema views}

We evaluate seven schema projections on the three trace corpora. The
LLM-focused projections are mapped from public schema families: OTel
GenAI semantic conventions (main branch f489c1832459), OpenInference
semantic conventions (main branch 0ee46c6c0d8e), and OpenLLMetry
v0.61.0 OTel application traces~\cite{otel_genai,openllmetry,openinference}.
We map published span-kind and attribute families into the local
normalized event vocabulary; a field is retained only when the cited
convention exposes the corresponding span kind or attribute family, or
when the projection is explicitly a custom-attribute or agent-structural
view. Fields with no public anchor are counted absent. This makes the
study a controlled field-retention analysis under the projection rules
in Table~\ref{tab:projections}, not a benchmark of every deployment of
those projects.

\begin{table*}[t]
\centering
\scriptsize
\caption{Schema projections evaluated by replaying sanitized events
through the OpenTelemetry SDK and a SQLite trace backend. Public anchors
are OTel GenAI conventions, OpenInference span kinds, and
OpenLLMetry-style OTel application traces; rows are field-family
mappings rather than product-deployment measurements.}
\label{tab:projections}
\begin{tabular}{@{}p{0.15\textwidth}p{0.19\textwidth}p{0.31\textwidth}p{0.27\textwidth}@{}}
\toprule
Projection & Public anchor & Retained field families & Absent field families \\
\midrule
Vanilla OTel & OTel trace model & trace/span/parent IDs, timestamps, status & GenAI and agent diagnostic attributes \\
OTel+custom attrs & OTel custom span attributes & \texttt{agent.event\_kind} and non-content diagnostic attributes & dedicated event-kind column and payload content \\
GenAI-only & OTel GenAI spans & \texttt{gen\_ai.*} LLM/tool fields plus retrieval evidence & planning, reasoning, guardrail, delegation, memory, task labels \\
LLM+agent-run & OpenInference span kinds & \texttt{AGENT}, \texttt{LLM}, \texttt{TOOL}, \texttt{RETRIEVER}, \texttt{GUARDRAIL} & planning, reasoning loops, delegation, memory, tool-argument shape \\
LLM app & OpenLLMetry-style app traces & agent, LLM, tool, retrieval, token, and tool-shape fields & planning, reasoning, guardrail, delegation, memory \\
Agent structural & agent structural extension & all event kinds plus structural and non-content diagnostic fields & prompt/completion/reasoning/retrieval/tool payload content \\
Agent reference & local reference schema & all sanitized fields needed by the local artifacts & raw payload content \\
\bottomrule
\end{tabular}
\end{table*}

\begin{table*}[t]
\centering
\scriptsize
\caption{SDK-to-SQLite backend validation and cost. The full generated
backend contains 740,272 span rows and 3,463,458 attribute rows across
projections and occupies 1178.9 MiB on disk. SQL ms is the combined time
for three backend-native aggregate queries.}
\label{tab:backend}
\begin{tabular}{@{}lrrrrr@{}}
\toprule
Projection & Spans & Attr rows & Export s & SQL ms & Eval s \\
\midrule
Vanilla OTel & 46267 & 0 & 1.15 & 0.5 & 0.84 \\
OTel+custom attrs & 46267 & 258617 & 1.95 & 21.7 & 1.25 \\
GenAI-only & 46267 & 169103 & 1.61 & 24.1 & 1.09 \\
LLM+agent-run & 46267 & 173799 & 1.53 & 23.9 & 1.06 \\
LLM app & 46267 & 187924 & 1.54 & 24.2 & 1.06 \\
Agent structural & 46267 & 258617 & 1.72 & 27.7 & 1.27 \\
Agent reference & 46267 & 259741 & 1.65 & 29.0 & 1.30 \\
\bottomrule
\end{tabular}
\end{table*}

\textbf{Vanilla OTel view.} This conservative degradation removes the
agent-specific analytical schema. It represents the case where spans may
exist in a tracing backend but the downstream analyst cannot query agent
span kinds or agent-specific fields.

\textbf{OTel+custom-attributes view.} This projection keeps the generic
OTel span envelope and non-content diagnostic custom attributes while
removing a dedicated event-kind column and payload content. It models a
production team that keeps custom attributes inside a generic tracing
backend.

\textbf{GenAI-only view.} This projection preserves LLM, tool, and
retrieval fields corresponding to GenAI-oriented semantic conventions.
It tests how far an LLM-centric schema can go without agent
orchestration events.

\textbf{LLM+agent-run view.} This projection preserves
LLM/tool/retrieval, agent-run, and guardrail fields exposed by
OpenInference-style span kinds. It does not preserve planning, memory,
delegation, or separate reasoning-loop event structure.

\textbf{LLM app view.} This view retains LLM calls, tool calls,
retrieval, and agent-level events while dropping planning, reasoning,
guardrail, delegation, and memory semantics. We use it as an
LLM-application trace schema-family mapping, not as a claim about every
deployment of any specific system.

\textbf{Agent structural view.} This view preserves agent-specific event
kinds and structural fields while redacting prompt, completion, tool
payload, retrieved-query, reasoning-chain, and task content.

\textbf{Agent reference view.} This view retains all sanitized fields
needed by the local artifacts and acts as the query-answerability
reference. Raw payload content is not written to the paper directory.

\subsection{Backend validation}

Each projection is replayed through the OpenTelemetry SDK and exported
to a SQLite trace-backend instance before metrics are computed. The
backend has separate span and attribute tables and materializes 740,272
span rows plus 3,463,458 attribute rows across schema views and
ablations. Representative aggregate queries are executed directly in
SQL, and the query-answerability metrics are computed from backend rows
rather than raw source records. This is not a benchmark of a production
tracing backend. It validates the data-engineering path studied here:
source events become spans, schema mappings determine visible fields, an
exporter writes a backend representation, and downstream analytics run
over that representation. The volume and velocity claims in this paper
are therefore limited to controlled backend materialization and query
execution, not production-scale ingestion.

As an additional backend-path check, an OTel Collector otelcol-contrib
version 0.153.0 \texttt{otlpjsonfile}$\rightarrow$\texttt{file} round
trip over the OTel+custom-attributes projection preserved 46,267/46,267
spans, 258,617/258,617 diagnostic span attributes, and 44,119/44,119
parent links; the file exporter added 46,267 bookkeeping attributes that
are excluded from the study queries.

For the controlled matrix, we use six telemetry conditions already
present in the result file: no telemetry, vanilla OTel, OTel GenAI,
OpenInference, structure-only agent traces, and full-capture agent
traces.

\subsection{Pipeline transformations}

The analysis treats schema loss as a transformation problem. A trace
pipeline can damage diagnostic quality in at least four ways.

\textbf{Kind erasure} maps semantically distinct events to a generic
category. A reasoning step, planning step, and memory read may all
survive as spans, but if they are indistinguishable \texttt{INTERNAL}
events, downstream queries cannot use them as agent semantics.

\textbf{Attribute stripping} removes fields while preserving event kind.
For example, a memory span without operation or key namespace may still
count as a memory event but cannot support memory-corruption analysis. A
guardrail span without outcome cannot support safety auditing.

\textbf{Join-key loss} breaks relations among spans. Agent analytics
often need to connect a planning step to later tool calls, a retrieval
event to the LLM answer it grounded, or a memory read to a later
reasoning span. Removing parent-child links, task identifiers, or agent
identifiers turns a structured trace into an event bag.

\textbf{Over-redaction} treats content removal as schema deletion.
Filtering should be able to remove prompt text, completion text,
retrieved query content, tool payloads, reasoning-chain text, or memory
values while retaining structural fields. If filtering deletes the span
kind or join keys, it can destroy measured structural value
unnecessarily.

\subsection{Fault families}

We separate fourteen fault classes into two families.
\emph{Operational} faults are detectable from generic execution, cost,
timeout, or tool fields in this corpus: context overflow, cost
explosion, infinite loop, timeout, tool failure, and wrong tool.
\emph{Semantic agent} faults require planning, reasoning, guardrail,
delegation, retrieval freshness, or memory fields: agent misroute,
circular delegation, guardrail bypass, hallucination, memory corruption,
planning failure, reasoning loop, and stale retrieval. This split is
intentionally diagnostic rather than taxonomic; it asks which fields are
needed by the analysis pipeline.

\subsection{Reproducibility}

All tables in this paper are regenerated by scripts in the paper's
\texttt{scripts/} directory. The default packaged path reads
\texttt{generated/sanitized\_events.jsonl},
\texttt{generated/tau\_truth\_labels.jsonl}, and
\texttt{generated/trace\_quality\_metrics.json}; it then replays schema
projections through the OpenTelemetry SDK, rebuilds the SQLite backend,
runs structural visibility queries, applies field ablations, refreshes
machine-readable metrics and LaTeX table snippets, and recompiles the
PDF. An optional \texttt{--from-raw} mode rebuilds the sanitized events
and summary metrics when the upstream result directories are available.
The generated artifact manifest lists the generated table snippets and
metric files, including the files used by the paper. Raw traces,
conversation text, and tool payloads are not copied into the submission
artifacts.

\begin{table}[t]
\centering
\scriptsize
\caption{Schema-interoperability metrics over the pooled backend. Macro
Q gives each positive structural query family equal weight; Inst. Q
pools query instances and is dominated by high-count token/tool
families.}
\label{tab:interop}
\begin{tabular}{@{}lrrrr@{}}
\toprule
Projection & Kinds & Attrs & Macro Q & Inst. Q \\
\midrule
Vanilla OTel & 0.0\% & 0.0\% & 0.0\% & 0.00\% \\
OTel+custom attrs & 0.0\% & 100.0\% & 100.0\% & 100.00\% \\
GenAI-only & 33.3\% & 31.6\% & 37.5\% & 99.97\% \\
LLM+agent-run & 55.6\% & 47.4\% & 50.0\% & 99.97\% \\
LLM app & 44.4\% & 47.4\% & 37.5\% & 99.97\% \\
Agent structural & 100.0\% & 100.0\% & 100.0\% & 100.00\% \\
Agent reference & 100.0\% & 100.0\% & 100.0\% & 100.00\% \\
\bottomrule
\end{tabular}
\end{table}

%% =====================================================================
\section{Results}
\label{sec:results}

\subsection{\textbf{RQ1}: LLM schemas preserve partial queryability}

\begin{table}[t]
\centering
\scriptsize
\caption{Ground-truth denominators for structural visibility queries in
the reference backend.}
\label{tab:denominators}
\begin{tabular}{@{}lr@{}}
\toprule
Query & Reference count \\
\midrule
Loops & 113 \\
Plan & 162 \\
Guard & 80 \\
Memory & 257 \\
Deleg. & 109 \\
Tool prov. & 14427 \\
Tokens & 2167286 \\
Retr. & 685 \\
\bottomrule
\end{tabular}
\end{table}

Table~\ref{tab:interop} reports the schema-interoperability metrics over
the pooled backend. The key result is not that generic schemas are
useless; it is that they are selectively useful. The GenAI-only,
LLM+agent-run, and LLM app projections preserve enough fields to support
tool, token, retrieval, and $\tau$-bench expected-action evidence.
However, they retain only 33.3--55.6\% of the semantic event kinds
observed in the source corpora and answer only 37.5--50.0\% of positive
structural visibility query families overall. The OTel+custom-attributes
projection is the counterexample: it has 0.0\% event-kind column
coverage, but it preserves 100.0\% structural query recall because
non-content diagnostic attributes remain queryable. The result is
therefore not ``OTel cannot work''; it is that the fields needed by
these structural queries need to survive somewhere in the backend
schema.

Metric choice is a visible sensitivity point rather than a hidden
detail. Macro Q gives each diagnostic query family equal weight. If
query instances are pooled, LLM-focused projections reach 99.97\% recall
because token-accounting contributes 2,167,286 of 2,183,119 positive
instances. Over only the tool, token, and retrieval families, the
LLM-focused projections score 100.0\%; over the five orchestration
families, GenAI-only and LLM app score 0.0\%, while the
OpenInference-style projection scores 20.0\% by retaining guardrail
spans. We therefore use Macro Q as the primary measure of
diagnostic-family support and report Inst. Q as the throughput-weighted
countermetric.

The distinction between raw trace availability and semantic
availability is important. A tracing backend can report that an
execution is fully traced while the exported schema prevents a
downstream job from grouping spans by reasoning step, identifying the
guardrail check that blocked execution, or joining memory reads to later
model outputs. The quality defect is not missing storage; it is missing
analytical structure.

Table~\ref{tab:queryrecall} makes the selectivity concrete. LLM-focused
projections recover tool provenance, token-accounting, and
retrieval-grounding queries; the OpenInference-style projection also
recovers guardrail audit visibility. They do not recover planning,
reasoning-loop, memory, or delegation queries. These failures are not
row-loss artifacts. They are answerability failures: the projected trace
table no longer contains the event kinds and attributes required by the
query.

\begin{table*}[t]
\centering
\scriptsize
\caption{Structural visibility-query recall by schema projection.
$\tau$-bench action and outcome evidence are reported separately in
Table~\ref{tab:tauoutcome}.}
\label{tab:queryrecall}
\resizebox{\textwidth}{!}{%
\begin{tabular}{@{}lrrrrrrrr@{}}
\toprule
Projection & Loops & Plan & Guard & Memory & Deleg. & Tool & Tokens & Retr. \\
\midrule
Vanilla OTel & 0.0\% & 0.0\% & 0.0\% & 0.0\% & 0.0\% & 0.0\% & 0.0\% & 0.0\% \\
OTel+custom attrs & 100.0\% & 100.0\% & 100.0\% & 100.0\% & 100.0\% & 100.0\% & 100.0\% & 100.0\% \\
GenAI-only & 0.0\% & 0.0\% & 0.0\% & 0.0\% & 0.0\% & 100.0\% & 100.0\% & 100.0\% \\
LLM+agent-run & 0.0\% & 0.0\% & 100.0\% & 0.0\% & 0.0\% & 100.0\% & 100.0\% & 100.0\% \\
LLM app & 0.0\% & 0.0\% & 0.0\% & 0.0\% & 0.0\% & 100.0\% & 100.0\% & 100.0\% \\
Agent structural & 100.0\% & 100.0\% & 100.0\% & 100.0\% & 100.0\% & 100.0\% & 100.0\% & 100.0\% \\
Agent reference & 100.0\% & 100.0\% & 100.0\% & 100.0\% & 100.0\% & 100.0\% & 100.0\% & 100.0\% \\
\bottomrule
\end{tabular}%
}
\end{table*}

\begin{table}[t]
\centering
\scriptsize
\caption{Per-corpus structural query recall. This separates the
independent $\tau$-bench workload from the smaller orchestration-heavy
trace corpora.}
\label{tab:percorpus}
\begin{tabular}{@{}lrrrr@{}}
\toprule
Corpus & Pos. queries & GenAI-only & LLM app & Agent structural \\
\midrule
SWE-bench & 7 & 42.9\% & 42.9\% & 100.0\% \\
E2E/topology & 8 & 37.5\% & 37.5\% & 100.0\% \\
$\tau$-bench & 2 & 100.0\% & 100.0\% & 100.0\% \\
\bottomrule
\end{tabular}
\end{table}

\begin{table}[t]
\centering
\scriptsize
\caption{$\tau$-bench outcome validation. Missing expected actions are
used as a schema-observable failure proxy and evaluated against external
task pass/fail labels. The final column is TP/FP/FN/TN for failure
prediction.}
\label{tab:tauoutcome}
\resizebox{\columnwidth}{!}{%
\begin{tabular}{@{}lrrrrrr@{}}
\toprule
Projection & Act. & Shape & Prec. & Recall & F1 & Conf. \\
\midrule
Vanilla OTel & 0.0\% & 0.0\% & 41.3\% & 97.0\% & 57.9\% & 773/1099/24/84 \\
OTel+custom attrs & 66.2\% & 66.2\% & 54.9\% & 61.4\% & 57.9\% & 489/402/308/781 \\
GenAI-only & 66.2\% & 0.0\% & 54.9\% & 61.4\% & 57.9\% & 489/402/308/781 \\
LLM+agent-run & 66.2\% & 0.0\% & 54.9\% & 61.4\% & 57.9\% & 489/402/308/781 \\
LLM app & 66.2\% & 66.2\% & 54.9\% & 61.4\% & 57.9\% & 489/402/308/781 \\
Agent structural & 66.2\% & 66.2\% & 54.9\% & 61.4\% & 57.9\% & 489/402/308/781 \\
Agent reference & 66.2\% & 66.2\% & 54.9\% & 61.4\% & 57.9\% & 489/402/308/781 \\
\bottomrule
\end{tabular}%
}
\end{table}

\begin{table}[t]
\centering
\scriptsize
\caption{$\tau$-bench pass/fail baselines for interpreting the
missing-action proxy. The final column is TP/FP/FN/TN.}
\label{tab:taubaselines}
\begin{tabular}{@{}lrrrrr@{}}
\toprule
Baseline & Acc. & Prec. & Recall & F1 & Conf. \\
\midrule
Always fail & 40.3\% & 40.3\% & 100.0\% & 57.4\% & 797/1183/0/0 \\
Always pass & 59.7\% & 0.0\% & 0.0\% & 0.0\% & 0/0/797/1183 \\
Tool-count mismatch & 46.0\% & 42.0\% & 90.2\% & 57.3\% & 719/992/78/191 \\
\bottomrule
\end{tabular}
\end{table}

\begin{table}[t]
\centering
\scriptsize
\caption{Per-domain/model $\tau$-bench failure-proxy validation for the
reference schema. The final column is TP/FP/FN/TN.}
\label{tab:taugroups}
\resizebox{\columnwidth}{!}{%
\begin{tabular}{@{}lrrrrrr@{}}
\toprule
Group & Act. & Shape & Prec. & Recall & F1 & Conf. \\
\midrule
gpt-4o airline & 67.1\% & 67.1\% & 75.9\% & 56.9\% & 65.0\% & 66/21/50/63 \\
gpt-4o retail & 62.8\% & 62.8\% & 54.7\% & 74.2\% & 62.9\% & 135/112/47/166 \\
sonnet-35 airline & 66.4\% & 66.4\% & 75.9\% & 46.8\% & 57.9\% & 101/32/115/152 \\
sonnet-35 retail & 67.7\% & 67.7\% & 44.1\% & 66.1\% & 52.9\% & 187/237/96/400 \\
\bottomrule
\end{tabular}%
}
\end{table}

Table~\ref{tab:percorpus} separates this result by corpus. The
$\tau$-bench workload is easier for LLM-focused schemas because it is
primarily a tool-calling workload. The smaller SWE-bench and
E2E/topology corpora are more orchestration-heavy, and their recall
drops to 37.5--42.9\% for the LLM-focused projections.

The $\tau$-bench outcome analysis is a negative control rather than a
claim of strong outcome diagnosis. Missing expected actions are only a
proxy for task failure: even with the reference schema, 5,880 of 8,880
expected actions are visible (66.2\% action and argument-shape recall)
and the proxy reaches 57.9\% F1 against task pass/fail labels, barely
above the always-fail F1 baseline. This deliberately weaker result shows
that schema preservation makes tool-sequence and argument-shape evidence
available, but task outcome diagnosis still needs domain semantics and
argument-value validation.

\subsection{\textbf{RQ2}: Individual field classes are load-bearing}

The field-ablation study in Table~\ref{tab:ablation} separates schema
design mistakes that a coarse projection table hides. Dropping both the
event-kind column and event-kind custom attribute collapses query recall
to 0.0\% even though other attributes remain. Collapsing only
orchestration event kinds drops query recall to 37.5\%, matching the
GenAI-only and LLM app projections in this study. Dropping only parent
identifiers preserves event kinds and attributes but breaks tool
provenance. Dropping only timestamps breaks reasoning-loop queries
because repeated reasoning is an ordered phenomenon. Dropping only
memory, guardrail, or token/cost attributes breaks the corresponding
query family.

The positive result is narrower but still important: content redaction
preserves 100.0\% recall for the structural visibility queries in this
workload. These queries require event kinds, ordering, join keys, and
selected non-content attributes; they do not require prompt, completion,
reasoning-chain text, tool payload, or retrieved query content. This
does not mean content is unnecessary for every diagnostic task.
Hallucination adjudication, stale-retrieval judgment, and argument-value
debugging can require payload inspection. The result instead identifies
the structural tasks that remain available under payload redaction, with
ordinary access controls and retention policies still required for
structural fields.

\begin{table*}[t]
\centering
\scriptsize
\caption{Field-ablation study on the reference schema. Each row removes
one field class while keeping the rest of the backend table intact.}
\label{tab:ablation}
\begin{tabular}{@{}lrrrrp{0.36\textwidth}@{}}
\toprule
Ablation & Kinds & Attrs & Joins & Query recall & Failed positive queries \\
\midrule
Drop semantic kind & 0.0\% & 94.7\% & 100.0\% & 0.0\% & reasoning\_loops, planning\_events, guardrail\_audit, memory\_ops, delegation\_edges, tool\_provenance, token\_accounting, retrieval\_grounding \\
Collapse orchestration & 44.4\% & 94.7\% & 100.0\% & 37.5\% & reasoning\_loops, planning\_events, guardrail\_audit, memory\_ops, delegation\_edges \\
Drop parent IDs & 100.0\% & 100.0\% & 0.0\% & 87.5\% & tool\_provenance \\
Drop timestamps & 100.0\% & 100.0\% & 100.0\% & 87.5\% & reasoning\_loops \\
Drop agent IDs & 100.0\% & 89.5\% & 100.0\% & 100.0\% & none \\
Drop memory attrs & 100.0\% & 89.5\% & 100.0\% & 87.5\% & memory\_ops \\
Drop guardrail attrs & 100.0\% & 94.7\% & 100.0\% & 87.5\% & guardrail\_audit \\
Drop token/cost & 100.0\% & 89.5\% & 100.0\% & 87.5\% & token\_accounting \\
Content redaction & 100.0\% & 100.0\% & 100.0\% & 100.0\% & none \\
\bottomrule
\end{tabular}
\end{table*}

\subsection{Validation against prior diagnostic proxies}

The diagnostic-proxy effect is measurable. Table~\ref{tab:diagnostic}
compares agent-semantics traces with a degraded OTel view on failed
SWE-bench traces. With agent semantics preserved, the mean root-cause
candidate set is 3.0 spans; under degradation it is 9.5 spans. The mean
number of spans inspected before the proxy diagnosis rises from 9.5 to
28.5. The trace still exists, but the missing schema forces the
automated diagnostic procedure to search a broader, less structured
execution space.

\begin{table}[t]
\centering
\caption{Downstream diagnostic degradation when agent-specific trace data
are removed. Lower is better for candidate spans and
spans-to-diagnosis; higher is better for signal-to-noise.}
\label{tab:diagnostic}
\resizebox{\columnwidth}{!}{%
\scriptsize
\begin{tabular}{@{}lrrr@{}}
\toprule
Metric & AgentTelemetry & Degraded OTel & Effect \\
\midrule
Candidate spans at root cause & 3.0 & 9.5 & 68.4\% fewer \\
Spans to diagnosis & 9.5 & 28.5 & 66.8\% fewer \\
Diagnostic signal-to-noise & 0.54 & 0.00 & semantic signal retained \\
\bottomrule
\end{tabular}
}
\end{table}

\begin{table}[t]
\centering
\scriptsize
\caption{Controlled fault-detection results split by whether the fault
requires agent-orchestration semantics. Operational faults can often be
detected from generic LLM/tool/cost traces; semantic faults degrade
sharply when planning, reasoning, guardrail, delegation, or memory
fields are absent.}
\label{tab:faultfamilies}
\begin{tabular}{@{}lrr@{}}
\toprule
Telemetry condition & Semantic faults & Operational faults \\
\midrule
No telemetry & 0.0\% & 0.0\% \\
Vanilla OTel & 0.0\% & 100.0\% \\
OTel GenAI & 0.0\% & 100.0\% \\
OpenInference & 0.0\% & 100.0\% \\
Structure-only & 32.1\% & 100.0\% \\
Full capture & 32.1\% & 100.0\% \\
\bottomrule
\end{tabular}
\end{table}

\begin{table}[t]
\centering
\scriptsize
\caption{Class-wise controlled semantic-fault detection. Preserving
schema fields is necessary but not sufficient: several classes also
require application assertions, expected-span specifications, or
argument/content checks.}
\label{tab:classwise}
\begin{tabular}{@{}lrrrr@{}}
\toprule
Fault class & OTel & OpenInf. & Struct. & Full \\
\midrule
agent\_misroute & 0.0\% & 0.0\% & 14.3\% & 14.3\% \\
circular\_delegation & 0.0\% & 0.0\% & 100.0\% & 100.0\% \\
guardrail\_bypass & 0.0\% & 0.0\% & 14.3\% & 14.3\% \\
hallucination & 0.0\% & 0.0\% & 71.4\% & 71.4\% \\
memory\_corruption & 0.0\% & 0.0\% & 14.3\% & 14.3\% \\
planning\_failure & 0.0\% & 0.0\% & 14.3\% & 14.3\% \\
reasoning\_loop & 0.0\% & 0.0\% & 14.3\% & 14.3\% \\
stale\_retrieval & 0.0\% & 0.0\% & 14.3\% & 14.3\% \\
\bottomrule
\end{tabular}
\end{table}

\begin{table}[t]
\centering
\scriptsize
\caption{Calibration sensitivity. Trace-data schemas that preserve agent
semantics still require threshold governance: raising retry or
token-growth thresholds reduces alert volume but can erase a diagnostic
class. Asterisks mark defaults.}
\label{tab:calibration}
\begin{tabular}{@{}lrrr@{}}
\toprule
Parameter & Value & SWE-bench alerts & Real-LLM alerts \\
\midrule
max\_retries & 2 & 15 & 34 \\
max\_retries & 3* & 7 & 26 \\
max\_retries & 5 & 3 & 6 \\
token\_growth\_factor & 1.15 & 25 & 59 \\
token\_growth\_factor & 1.3* & 7 & 26 \\
token\_growth\_factor & 1.45 & 4 & 21 \\
token\_growth\_factor & 2.0 & 4 & 21 \\
\bottomrule
\end{tabular}
\end{table}

This is the practical meaning of trace-data quality for the measured
LLM-agent diagnostics. Schema loss does not merely remove descriptive
labels. It changes the cost of downstream structural analytics. A
reasoning-loop detector cannot fire if reasoning spans are not
represented. A guardrail audit cannot count blocked safety events if
guardrail spans are flattened away. A memory-corruption analysis cannot
connect a stale memory read to a later action if memory fields are
absent.

The controlled fault matrix gives a second view of the same effect.
Table~\ref{tab:faultfamilies} separates operational faults from semantic
agent faults. Generic telemetry detects operational failures in this
corpus because those faults are exposed through generic execution
signals such as timeout, cost, or tool error. It detects 0.0\% of
semantic agent faults because the required fields are absent.
Structure-only and full-capture agent traces recover a subset of
semantic-fault detections while retaining 100.0\% operational-fault
detection. The incomplete semantic detection rate is itself useful:
schema preservation is necessary but not sufficient. Some agent faults
also require framework support, expected-span specifications, or
application-specific assertions. Table~\ref{tab:classwise} shows the
class-wise pattern: circular delegation is almost entirely structural in
this matrix, while most semantic classes need more than event-kind
preservation.

This separation is useful for system design. If the only goal is to
catch timeouts, tool exceptions, or cost spikes, generic LLM/tool
telemetry may be enough. If the goal is to debug why an agent selected
the wrong action, entered a reasoning loop, bypassed a guardrail, or
used stale memory, generic telemetry has the wrong data model. A data
platform should not advertise those diagnostic tasks unless it retains
the corresponding schema fields.

\subsection{\textbf{RQ3}: Calibration remains a separate risk}

However, schema quality does not remove calibration risk.
Table~\ref{tab:calibration} shows threshold sensitivity for selected
detector parameters. Raising retry or token-growth thresholds reduces
alert volume, but it can also erase a diagnostic class. For example,
increasing the token-growth threshold from the default 1.3 to 1.45
reduces real-LLM alerts from 26 to 21 in the underlying result file.
That reduction is not automatically a quality improvement; it may be a
false-negative tradeoff. Trace-data quality therefore has two layers:
the schema needs to preserve the fields needed by the target diagnostic
query, and the analytics must calibrate thresholds without silently
disabling the failure modes the schema made visible.

%% =====================================================================
\section{Discussion}
\label{sec:discussion}

\subsection{Agent observability as data engineering}

The results suggest that LLM-agent observability should be designed as
a data pipeline, not just as an SDK feature. A production trace system
needs a schema contract for agent events, collector transformations that
do not erase orchestration semantics, storage schemas that retain join
keys, and analytics that distinguish content capture from
structural-field capture. Without those properties, a large trace lake
can still be analytically low quality for agent diagnostics.

This framing also explains why LLM-only observability is insufficient
for agents. LLM calls are important rows in the trace table, but they
are not the whole table. Agent diagnosis often depends on the rows
around the LLM call: the plan that chose it, the reasoning step that
interpreted it, the retrieval event that grounded it, the memory record
that biased it, and the guardrail that approved or blocked it.

\subsection{Payload-redacted structural visibility}

The payload-redacted structural result is encouraging for a bounded
class of diagnostics. Organizations do not always need to capture
prompt, completion, retrieved-query, reasoning-chain, or tool-payload
content to answer structural visibility questions about loops,
guardrails, memory operations, delegation, token accounting, and tool
provenance. Payload content remains necessary for other tasks, such as
judging whether a retrieved document actually grounded an answer or
whether a tool argument had the correct value. This does not eliminate
sensitivity, because even structural fields can expose behavior and
workflow shape, but it creates a better engineering tradeoff than
all-or-nothing trace capture.

\subsection{Why BigData is the right framing}

The same empirical facts could be misframed as a tool comparison or an
agent-debugging result. That would understate the broader systems
lesson. The IEEE BigData framing is about the 5Vs of data-intensive
systems. LLM-agent traces create volume in the form of many structured
events, variety through heterogeneous producer schemas, velocity as
traces move through collectors and backends, value through downstream
diagnostic analytics, and veracity risks when schema transformations
drop the fields that make trace rows meaningful. This paper demonstrates
those risks on a controlled, backend-materialized corpus rather than a
production ingestion benchmark. It therefore belongs with data and
information quality, big-data infrastructure, and data-processing
pipelines, not only with agent evaluation.

\subsection{Implications for schema design}

A useful agent-trace schema should satisfy four constraints. First, it
should preserve agent-orchestration event kinds separately from generic
internal spans. Second, it should make structural joins cheap by
carrying trace, task, agent, parent, and session identifiers where
appropriate. Third, it should allow privacy filters to redact payload
content without deleting structural fields. Fourth, it should document
the diagnostic tasks it supports. A schema that cannot answer whether a
reasoning loop occurred or whether a guardrail fired should not be
treated as complete agent observability.

\subsection{Open implementation questions}

Several engineering questions remain open. First, agent-trace schemas
need versioning. A collector or warehouse should know whether a trace
uses a generic, GenAI-only, LLM+agent-run, LLM app, or agent-specific
extension, and should record which semantic fields were intentionally
unavailable. Second, schema transformations should expose loss reports:
if a processor drops reasoning or memory fields, that should be
measurable provenance rather than an invisible side effect. Third,
privacy filters should separate content sensitivity from structural
sensitivity, allowing operators to keep diagnostic structure while
redacting payloads. Finally, evaluations for observability systems
should include downstream queries, not only ingestion throughput or
export success.

%% =====================================================================
\section{Threats to Validity}
\label{sec:threats}

\textbf{Schema degradation is an analytical model.} The generic and
LLM-focused views are controlled field-retention mappings from public
schema families, not every possible deployment of a named system. A
production team could add custom attributes to generic spans and recover
some signal. That would support, rather than weaken, the paper's thesis:
the data-quality outcome depends on which semantic fields are preserved.

\textbf{The backend evidence is not a production benchmark.} The
OpenTelemetry SDK-to-SQLite path validates export, storage, and SQL
querying over projected spans. The Collector round trip checks that one
receiver/exporter path preserves spans, diagnostic attributes, and
parent links for the custom-attribute projection, but it does not
measure Jaeger, Tempo, ClickHouse, or search-backend behavior. The
backend numbers should be read as trace-lake cost evidence for this
study, not as production observability-system performance results.

\textbf{The corpora are still limited.} Two trace corpora and the
controlled matrix come from the \textsc{AgentTelemetry} artifact, while
the $\tau$-bench trajectory corpus is independent and uses external task
labels. This improves independence but does not cover every agent
framework, domain, or observability backend. Future work should repeat
the same schema-loss analysis on additional production trace corpora
and backend deployments.

\textbf{Semantic faults are harder than operational faults.} The
controlled matrix shows that even full-capture and structure-only agent
traces do not detect all semantic faults in all frameworks. This is not
hidden in the analysis. It means schema preservation is a necessary
condition for semantic fault diagnosis, not a complete solution.

\textbf{Diagnostic metrics are proxies.} Candidate spans, spans to
diagnosis, and signal-to-noise are automated proxies for diagnostic
effort. They do not replace a human operator study. They are appropriate
for this paper because the question is data-quality degradation in the
trace pipeline, not human incident response performance.

\textbf{The paper uses prior artifacts.} The paper intentionally builds
on an accepted prior artifact. The risk is perceived overlap. We
mitigate that by citing the prior work as the source of the span/fault
definitions, avoiding benchmark or toolkit novelty claims, and limiting
the new claims to schema-loss, data-quality, and diagnostic-degradation
analysis.

%% =====================================================================
\section{Related Work}
\label{sec:related}

\textbf{Distributed tracing and observability.} Classical tracing
systems such as Dapper, X-Trace, and Magpie established the value of
structured request traces for large-scale
systems~\cite{dapper,xtrace,magpie}. OTel generalized these ideas into
vendor-neutral APIs, SDKs, collectors, and semantic conventions for
modern production systems~\cite{otel_spec}. This paper studies a new
schema-quality problem that arises when the traced computation is an
LLM-agent workflow rather than a service call graph.

\textbf{Data quality, schema interoperability, and provenance.}
Data-quality research distinguishes syntactic correctness from fitness
for consumer tasks~\cite{wang_strong,batini}. Schema-matching and
data-integration work studies how heterogeneous schemas can be aligned
and where mappings fail~\cite{rahm_bernstein}. Database provenance work
shows why downstream query answers depend on which source records and
transformations survive the pipeline~\cite{cheney_provenance}. This
paper applies those ideas to observability traces: a span table is not
high quality merely because rows are present; it is high quality only if
it preserves the fields needed by diagnostic consumers.

\textbf{LLM and agent observability.} OTel GenAI, OpenInference,
OpenLLMetry, LangSmith, and Langfuse all address parts of LLM
observability~\cite{otel_genai,openllmetry,langsmith,langfuse,openinference}.
Their schemas make LLM calls, tool calls, retrieval, and application
traces easier to analyze. Our analysis shows why agent-level
orchestration semantics are needed by the measured diagnostic workloads.

\textbf{Agent failure analysis.} Recent work documents multi-agent and
LLM-agent failure modes, debugging workflows, and observability
needs~\cite{mast,agentops,agdebugger}. The prior \textsc{AgentTelemetry}
paper introduced a benchmark and toolkit for agent fault
detection~\cite{agenttelemetry_aiware}. $\tau$-bench provides a
complementary benchmark for tool-agent-user interactions with task
outcomes and expected actions~\cite{taubench}. This paper studies the
data-quality and schema-loss properties of telemetry pipelines that
carry such traces.

%% =====================================================================
\section{Conclusion}
\label{sec:conclusion}

LLM-agent observability is a big-data quality problem. A pipeline can
store thousands of spans while losing the semantic structure needed for
diagnosis. Across 2,148 traces and 46,267 backend-exported events,
LLM-focused schema projections preserve tool, token, and retrieval query
families but answer only 37.5--50.0\% of positive structural visibility
query families overall and 0.0\% of the measured planning,
reasoning-loop, memory, and delegation families. Field ablations show
that event kind, parent identifiers, timestamps, memory fields,
guardrail outcomes, and token/cost fields each carry distinct diagnostic
load. A payload-redacted structural projection preserves all measured
structural visibility answers without retaining prompt, completion,
reasoning-chain, retrieved-query, or tool payload content. A generic
OTel span store can also preserve these structural queries when
non-content custom diagnostic attributes survive export. Agent traces
therefore need explicit schema-quality contracts: if planning,
reasoning, guardrail, delegation, and memory fields do not survive
ingestion and transformation, analytics can run on complete-looking
traces while answering incomplete questions.

%% =====================================================================
%% References: hand-written in the exact form IEEEtran.bst (v1.14)
%% produced for the submitted PDF.
\balance
\begin{thebibliography}{10}
\providecommand{\url}[1]{#1}
\csname url@samestyle\endcsname
\providecommand{\newblock}{\relax}
\providecommand{\bibinfo}[2]{#2}
\providecommand{\BIBentrySTDinterwordspacing}{\spaceskip=0pt\relax}
\providecommand{\BIBentryALTinterwordstretchfactor}{4}
\providecommand{\BIBentryALTinterwordspacing}{\spaceskip=\fontdimen2\font plus
\BIBentryALTinterwordstretchfactor\fontdimen3\font minus
  \fontdimen4\font\relax}
\providecommand{\BIBforeignlanguage}[2]{{%
\expandafter\ifx\csname l@#1\endcsname\relax
\typeout{** WARNING: IEEEtran.bst: No hyphenation pattern has been}%
\typeout{** loaded for the language `#1'. Using the pattern for}%
\typeout{** the default language instead.}%
\else
\language=\csname l@#1\endcsname
\fi
#2}}
\providecommand{\BIBdecl}{\relax}
\BIBdecl

\bibitem{otel_spec}
\BIBentryALTinterwordspacing
{OpenTelemetry Authors}, ``{OpenTelemetry Specification},'' 2024. [Online].
  Available: \url{https://opentelemetry.io/docs/specs/otel/}
\BIBentrySTDinterwordspacing

\bibitem{otel_genai}
------, ``{OpenTelemetry} semantic conventions for {GenAI} operations,''
  \url{https://opentelemetry.io/docs/specs/semconv/gen-ai/}, 2025,
  f489c1832459, Jun. 2026.

\bibitem{openllmetry}
{Traceloop}, ``{OpenLLMetry}: Open-source observability for {LLM}
  applications,'' \url{https://github.com/traceloop/openllmetry}, 2025,
  v0.61.0, Jun. 2026.

\bibitem{langsmith}
{LangChain, Inc.}, ``{LangSmith}: Platform for {LLM} application development,
  monitoring, and testing,'' \url{https://smith.langchain.com}, 2025.

\bibitem{langfuse}
{Langfuse}, ``{Langfuse}: Open-source {LLM} engineering platform,''
  \url{https://github.com/langfuse/langfuse}, 2025.

\bibitem{openinference}
\BIBentryALTinterwordspacing
{Arize AI}, ``{OpenInference}: Semantic conventions for {LLM}, ml, and agent
  observability,'' 2024, 0ee46c6c0d8e, Jun. 2026. [Online]. Available:
  \url{https://github.com/Arize-ai/openinference}
\BIBentrySTDinterwordspacing

\bibitem{agenttelemetry_aiware}
K.~C. Balusu, ``{AgentTelemetry}: A fault detection benchmark and toolkit for
  {LLM} agent observability,'' in \emph{Proceedings of the 3rd ACM
  International Conference on AI-Powered Software (AIware '26)}.\hskip 1em plus
  0.5em minus 0.4em\relax ACM, 2026, montreal, QC, Canada, July 6--7, 2026.

\bibitem{agenttelemetry_artifact}
------, ``{AgentTelemetry}: Reproduction artifact,'' 2026.

\bibitem{swebench}
\BIBentryALTinterwordspacing
C.~E. Jimenez, J.~Yang, A.~Wettig, S.~Yao, K.~Pei, O.~Press, and K.~Narasimhan,
  ``{SWE-bench}: Can language models resolve real-world {GitHub} issues?'' in
  \emph{Proceedings of the 12th International Conference on Learning
  Representations}, 2024. [Online]. Available:
  \url{https://openreview.net/forum?id=VTF8yNQM66}
\BIBentrySTDinterwordspacing

\bibitem{taubench}
\BIBentryALTinterwordspacing
S.~Yao, N.~Shinn, P.~Razavi, and K.~Narasimhan, ``$\tau$-bench: A benchmark for
  tool-agent-user interaction in real-world domains,'' 2024. [Online].
  Available: \url{https://arxiv.org/abs/2406.12045}
\BIBentrySTDinterwordspacing

\bibitem{dapper}
\BIBentryALTinterwordspacing
B.~H. Sigelman, L.~A. Barroso, M.~Burrows, P.~Stephenson, M.~Plakal, D.~Beaver,
  S.~Jaspan, and C.~Shanbhag, ``{Dapper, a Large-Scale Distributed Systems
  Tracing Infrastructure},'' Google, Tech. Rep. dapper-2010-1, 2010. [Online].
  Available:
  \url{https://research.google/pubs/dapper-a-large-scale-distributed-systems-tracing-infrastructure/}
\BIBentrySTDinterwordspacing

\bibitem{xtrace}
\BIBentryALTinterwordspacing
R.~Fonseca, G.~Porter, R.~H. Katz, S.~Shenker, and I.~Stoica, ``{X-Trace}: A
  pervasive network tracing framework,'' in \emph{Proceedings of the 4th USENIX
  Symposium on Networked Systems Design and Implementation}, 2007. [Online].
  Available: \url{https://www.usenix.org/legacy/event/nsdi07/tech/fonseca.html}
\BIBentrySTDinterwordspacing

\bibitem{magpie}
\BIBentryALTinterwordspacing
P.~Barham, A.~Donnelly, R.~Isaacs, and R.~Mortier, ``Using magpie for request
  extraction and workload modelling,'' in \emph{Proceedings of the 6th
  Symposium on Operating Systems Design and Implementation}, 2004. [Online].
  Available: \url{https://www.usenix.org/legacy/event/osdi04/tech/barham.html}
\BIBentrySTDinterwordspacing

\bibitem{mast}
M.~Cemri, M.~Z. Pan, S.~Yang, L.~A. Agrawal, B.~Chopra, R.~Tiwari, K.~Keutzer,
  A.~Parameswaran, D.~Klein, K.~Ramchandran, M.~Zaharia, J.~E. Gonzalez, and
  I.~Stoica, ``Why do multi-agent {LLM} systems fail?'' \emph{arXiv preprint
  arXiv:2503.13657}, 2025.

\bibitem{agentops}
L.~Dong, Q.~Lu, and L.~Zhu, ``{AgentOps}: Enabling observability of {LLM}
  agents,'' 2024.

\bibitem{agdebugger}
W.~Epperson, G.~Bansal, V.~Dibia, A.~Fourney, J.~Gerrits, E.~Zhu, and
  S.~Amershi, ``Interactive debugging and steering of multi-agent {AI}
  systems,'' in \emph{Proceedings of the 2025 CHI Conference on Human Factors
  in Computing Systems}, 2025.

\bibitem{wang_strong}
R.~Y. Wang and D.~M. Strong, ``Beyond accuracy: What data quality means to data
  consumers,'' \emph{Journal of Management Information Systems}, vol.~12,
  no.~4, pp. 5--33, 1996.

\bibitem{batini}
C.~Batini, C.~Cappiello, C.~Francalanci, and A.~Maurino, ``Methodologies for
  data quality assessment and improvement,'' \emph{ACM Computing Surveys},
  vol.~41, no.~3, pp. 1--52, 2009.

\bibitem{rahm_bernstein}
E.~Rahm and P.~A. Bernstein, ``A survey of approaches to automatic schema
  matching,'' \emph{The VLDB Journal}, vol.~10, no.~4, pp. 334--350, 2001.

\bibitem{cheney_provenance}
J.~Cheney, L.~Chiticariu, and W.-C. Tan, ``Provenance in databases: Why, how,
  and where,'' \emph{Foundations and Trends in Databases}, vol.~1, no.~4, pp.
  379--474, 2009.

\end{thebibliography}

\end{document}
